\chapter{Introduction}
\label{chap:introduction}

\section{Context}

Small and independent retailers, such as a local craft shop or a neighbourhood restaurant, face a structural disadvantage when it comes to marketing. Larger competitors employ in-house design teams and external creative agencies that produce a steady stream of polished promotional material throughout the year. Independent shops, by contrast, rarely have the time or the expertise to produce the kind of professional visual content that today's customers have come to expect on social media and in physical storefronts. The result is a visible quality gap between large brands and the local businesses that compete with them, and this gap continues to widen as larger competitors invest aggressively in digital marketing.

In the last few years, generative artificial intelligence has begun to close part of this gap. Modern image-generating models can produce near-photographic visuals from a short text prompt, and conversational language models can draft promotional copy on demand. In principle, these tools should be transformative for small retailers; in practice, however, the leap from a general-purpose model to a polished, brand-consistent advertisement is far from trivial. Models often produce renderings that look good in isolation but misrepresent the product itself, struggle to reproduce a retailer's exact branding (a particular logo or a chosen colour palette, for instance), and require careful prompting that most non-technical users are not equipped to perform. They also do not, on their own, decide when a campaign should be run, or what theme would resonate with a particular local audience on a particular day.

The platform presented in this thesis was developed against this backdrop. It treats marketing automation as more than a single image-generation step. It records the retailer's brand profile alongside a library of their products, watches a small set of real-world contextual signals (current weather, upcoming holidays, recent news) and uses them to suggest daily marketing ideas. When the retailer likes an idea, the platform can then turn it into a finished asset to publish. The client side runs on both desktop web browsers and mobile devices from a single React Native codebase, and the service is hosted on the Microsoft Azure cloud as a regular production deployment rather than as a controlled experiment.

\section{Objectives}
\label{sec:objectives}

Concretely, this thesis pursues the following objectives:

\begin{itemize}
    \item To design and implement a cross-platform application that automates the production of marketing assets directly from a retailer's product photographs and shop information.
    \item To enable retailers to assemble their product library with minimal manual effort, either by uploading their own photographs and removing the background automatically through an AI model, or by searching a curated reference database by visual similarity to a sample photo or by name.
    \item To develop a recommendation engine that suggests fresh marketing ideas on a daily basis, grounded in real-world contextual signals such as the weather, upcoming holidays and recent news, so that the platform is useful as an everyday tool rather than a one-off generator.
    \item To investigate and implement a hybrid generation strategy in which a generative model designs the surrounding scene while the retailer's actual product photograph is preserved through deterministic post-processing, addressing a well-known limitation of current models in faithfully reproducing specific real-world objects. The same pipeline is then extended with a high-resolution rendering path that produces print-ready posters for display in the shop itself.
    \item To package the resulting platform as a cloud service deployed on Microsoft Azure, to demonstrate that the proposed approach is viable beyond a controlled experimental setting.
\end{itemize}

\section{Thesis Structure}

The remainder of this thesis is organised as follows. Chapter 2 reviews the theoretical foundations on which the system is built: modern image-generating models, prompt engineering, retrieval-augmented generation, and techniques for measuring visual similarity between images. Chapter 3 is devoted to the overall architecture of the platform and to its deployment on the Microsoft Azure cloud. The implementation is then presented in Chapter 4, with a particular focus on the recommendation engine and the hybrid generation technique that together constitute the main technical contributions of this work. Chapter 5 covers the testing strategy, including both the automated suites for architectural correctness and the offline evaluation protocol applied to the language-model output. Chapter 6 closes the thesis with a summary of the contributions and a number of directions for future work.

During the preparation of this work the author used ChatGPT to improve the language and readability of the thesis. After using ChatGPT, I reviewed and edited the content as needed. All experimental results, model training, system design, evaluation, and scientific conclusions are entirely my own work. I take full responsibility for the accuracy and integrity of the final document.
